\newpage
\heading{Appendix B. Why the degree is exactly $h$}
{Lai, section 6, Proposition 6.1. Optional for the irrationality argument.}

The main proof only needs $\deg\Delta_K\leq h$. We can also show that
the highest-degree term does not cancel. The reason is that cancellation
in $A_K=D_N^6/D_K$ leaves exactly $h=K-N$ distinct poles, matching the
number of rows of the matrix.

Write $a_k=-k^2$ for $N<k\leq K$ and
$T_K(t)=\prod_{k=N+1}^K(t-a_k)$. The residue of
$A_K(t)t^{i+j}=D_N(t)^5t^{i+j}/T_K(t)$ at $a_k$ is
\[
 c_k^{ij}=\frac{D_N(a_k)^5a_k^{i+j}}{T_K'(a_k)}.
\]
Indeed, multiply by $t-a_k$ and evaluate at $a_k$; the other denominator
factors multiply to $T_K'(a_k)$. The polynomial part of the partial
fraction decomposition contributes no $X$ to $\mu_X$. By
\eqref{eq:mu-pole}, the coefficient of $X$ contributed by this pole
is $k^4c_k^{ij}$. Thus, writing $G_K(X)=G_0+XG_1$, we have
\[
 (G_1)_{ij}=\sum_{k=N+1}^K d_ka_k^ia_k^j,
 \qquad d_k=\frac{k^4D_N(a_k)^5}{T_K'(a_k)}.
\]

Let $V$ be the square matrix with entries $V_{ik}=a_k^i$, where
$0\leq i<h$ and $N<k\leq K$. The preceding equality is the matrix identity
\[
 G_1=V\operatorname{diag}(d_k)V^{\mathsf T},\qquad
 \det V=\prod_{N<k<\ell\leq K}(a_\ell-a_k).
\]
For completeness, the Vandermonde formula follows because the determinant
is an alternating polynomial in the nodes: it vanishes whenever two nodes
coincide, so every difference $a_\ell-a_k$ divides it. The product of these
differences has the same total degree $0+1+\cdots+(h-1)$. Comparing the
coefficient of the monomial whose successive exponents are
$0,1,\ldots,h-1$ shows that the remaining constant is one.

All nodes are distinct, so $\det V\ne0$. Also $d_k\ne0$: the factors
of $D_N(a_k)=\prod_{j=1}^N(j^2-k^2)$ are nonzero, as are those of
$T_K'(a_k)$. Therefore $\det G_1\ne0$. In the determinant expansion,
the coefficient of $X^h$ is obtained by taking the $X$ term from each of
the $h$ factors. Consequently
\[
 [X^h]\Delta_K(X)=\det G_1\ne0,\qquad \deg\Delta_K=h.
\]

The leading coefficient has an especially simple form. Pairing the two
factors belonging to each unordered pair of nodes gives
\[
 \prod_{k=N+1}^K T_K'(a_k)
 =(-1)^{h(h-1)/2}\prod_{N<k<\ell\leq K}(a_\ell-a_k)^2.
\]
Using this in $\det G_1=(\det V)^2\prod_k d_k$ cancels the squared
differences and yields
\[
 \boxed{[X^h]\Delta_K(X)
 =(-1)^{h(h-1)/2}\prod_{k=N+1}^K k^4D_N(-k^2)^5
 \in\Z\setminus\{0\}.}
\]
This concerns one coefficient. It does not assert that all coefficients
of $\Delta_K$ are integers, and it is separate from the positivity of
$\Delta_K(\xi)$.

\newpage
\heading{Appendix C. The same functional in different coordinates}
{Lai, section 7, Lemma 7.3. Optional once the arithmetic estimate is assumed.}

The auxiliary functional in section 7 changes the pole locations from
negative squares to integers. Here is the complete identity, including
the weight that accompanies the substitution.

Define $L:\Q[x]\to\Q$ by $L(x^k)=B_k$, and put
\[
 \tau(P)=\frac{L(P''')}{24},\qquad
 d(r)=\begin{cases}r,&r\geq0,\\-r-1,&r<0.\end{cases}
\]
On rational functions with only simple integer poles, define $\tau_X$
by its polynomial restriction $\tau$ and the pole values
\[
 \tau_X\!\left(\frac1{x-r}\right)=\Hfive{d(r)}-X,
 \qquad r\in\Z.
\]
Unique partial fractions extend these prescriptions linearly, just as
in the definition of $\mu_X$.

For every rational function $R$ in the domain of $\mu_X$, the identity is
\[
 \boxed{\mu_X(R)=\tau_X\!\left(x^5R(-x^2)\right).}
\]
The right side is defined: a pole at $t=-j^2$ becomes a pair of simple
poles at $x=j$ and $x=-j$, and a polynomial remains a polynomial after
substitution and multiplication by $x^5$. Both sides are linear in $R$.
By \eqref{eq:partial}, it is enough to check monomials and simple poles.

For $R(t)=t^e$, the transformed function is $(-1)^ex^{2e+5}$. Three
differentiations give
\[
 \tau_X\bigl((-1)^ex^{2e+5}\bigr)
 =\frac{(-1)^e(2e+5)(2e+4)(2e+3)B_{2e+2}}{24},
\]
which is exactly \eqref{eq:mu-moment}. This explains the order of the
derivative: it produces both the three consecutive factors and the
required Bernoulli index.

For $R(t)=1/(t+j^2)$, polynomial division gives
\[
 \frac{x^5}{j^2-x^2}
 =-x^3-j^2x-\frac{j^4}{2}
       \left(\frac1{x-j}+\frac1{x+j}\right).
\]
One can check the signs by multiplying by $j^2-x^2$. The polynomial
terms give $x^5-j^4x$, and the pole terms give $+j^4x$.
Now $\tau(-x^3)=-1/4$, $\tau(-j^2x)=0$, $d(j)=j$, and $d(-j)=j-1$.
Applying $\tau_X$ therefore gives
\begin{align*}
 \tau_X\!\left(\frac{x^5}{j^2-x^2}\right)
 &=-\frac14-\frac{j^4}{2}
       \left(\Hfive{j}+\Hfive{j-1}-2X\right)\\
 &=j^4\bigl(X-\Hfive{j}\bigr)-\frac14+\frac1{2j},
\end{align*}
where the last equality uses
$\Hfive{j}-\Hfive{j-1}=j^{-5}$. This is \eqref{eq:mu-pole}, completing
the proof by linearity.

Both the substitution $t=-x^2$ and the factor $x^5$ are necessary.
The identity supplies a second description of the same functional;
it introduces no new determinant. Integer pole locations make the
arithmetic calculations by residue class possible, while the original
$t$-description makes the positive integral representation visible.
